% ── Michigan TAC — PDF preamble ───────────────────────────────────────────
% Mirrors styles/_tac-core.scss so the two output channels read as one system.
% Colour values here are the light-theme tokens from styles/tac-light.scss.
% Engine: XeLaTeX (fontspec + unicode-math require it).

\RequirePackage{fontspec}
\RequirePackage{unicode-math}
\RequirePackage[svgnames]{xcolor}
\RequirePackage{tcolorbox}
\tcbuselibrary{most,skins,breakable}
\RequirePackage{caption}
\RequirePackage{enumitem}
\RequirePackage{booktabs}
\RequirePackage{longtable}
\RequirePackage{array}
\RequirePackage{ragged2e}
\RequirePackage{etoolbox}

% ── Fonts ─────────────────────────────────────────────────────────────────
% Resolved by family name against system fonts. CI must install Noto and
% instance Josefin Sans from its variable font — see .github/workflows/publish.yml.
\setmainfont{Noto Sans}[Scale=1.0, Ligatures=TeX]
\setsansfont{Noto Sans}[Scale=1.0, Ligatures=TeX]
% TeX text-ligatures deliberately disabled on mono: in verbatim they turn " into ”
% and -- into an en-dash, and the ToUnicode table records the substitution, so
% copy-pasted code no longer parses.
\setmonofont{Noto Sans Mono}[Scale=0.92, RawFeature={-tlig}]
\setmathfont{Noto Sans Math}[Scale=1.0, math-style=ISO, bold-style=ISO]

% Ligatures=TeX is required here too, not just on the body font: without it the
% `` and '' that Pandoc emits for smart quotes stay as literal backticks in any
% heading set in this family.
\newfontfamily\HeadingFont[Scale=MatchLowercase, Ligatures=TeX,
  BoldFont=Josefin Sans SemiBold]{Josefin Sans}
\newfontfamily\HeadingFontSemiBold[Scale=MatchLowercase, Ligatures=TeX]{Josefin Sans SemiBold}
\newfontfamily\HeadingFontLight[Scale=MatchLowercase, Ligatures=TeX]{Josefin Sans Light}

% ── Colour tokens (mirror of tac-light.scss :root) ────────────────────────
\definecolor{SurfacePage}{HTML}{EDEDEB}
\definecolor{SurfaceElevated}{HTML}{F7F7F6}
\definecolor{SurfaceCode}{HTML}{FAF2E3}
\definecolor{SurfaceSubtle}{HTML}{E9E9E7}
\definecolor{ContentPrimary}{HTML}{14140F}
\definecolor{ContentSecondary}{HTML}{424240}
\definecolor{ContentMuted}{HTML}{706964}
\definecolor{BorderHairline}{HTML}{DCDCDA}
\definecolor{BorderModerate}{HTML}{C8C8C5}
\definecolor{InteractivePrimary}{HTML}{434739}
\definecolor{InteractiveAccent}{HTML}{1F618D}
\definecolor{VerticalAccent}{HTML}{5A6310}
\definecolor{EditorialRule}{HTML}{EA8441}
\definecolor{StatusOpen}{HTML}{1F618D}
\definecolor{StatusAccepted}{HTML}{4A7C2F}
\definecolor{StatusDeclined}{HTML}{A33A2A}
\definecolor{StatusProgress}{HTML}{B06A10}

\pagecolor{SurfacePage}
\color{ContentSecondary}

% ── Headings in the display face ──────────────────────────────────────────
\usepackage{sectsty}
\allsectionsfont{\HeadingFont\color{ContentPrimary}}

% ── Quarto callouts, recoloured to the token set ──────────────────────────
% Quarto emits its own tcolorbox definitions; these redefine the frame colours
% so the PDF matches the SCSS rules in _tac-core.scss.
\definecolor{quarto-callout-color}{HTML}{1F618D}
\definecolor{quarto-callout-note-color}{HTML}{1F618D}
\definecolor{quarto-callout-tip-color}{HTML}{4A7C2F}
\definecolor{quarto-callout-important-color}{HTML}{A33A2A}
\definecolor{quarto-callout-warning-color}{HTML}{B06A10}
\definecolor{quarto-callout-caution-color}{HTML}{EA8441}
\definecolor{quarto-callout-note-color-frame}{HTML}{1F618D}
\definecolor{quarto-callout-tip-color-frame}{HTML}{4A7C2F}
\definecolor{quarto-callout-important-color-frame}{HTML}{A33A2A}
\definecolor{quarto-callout-warning-color-frame}{HTML}{B06A10}
\definecolor{quarto-callout-caution-color-frame}{HTML}{EA8441}

% ── Recommendation block — the PDF twin of .tac-recommendation ────────────
\newtcolorbox{tacrec}[1][]{
  enhanced, breakable,
  colback=SurfaceSubtle, colframe=InteractiveAccent,
  boxrule=0pt, leftrule=3pt, arc=2pt,
  left=9pt, right=9pt, top=7pt, bottom=7pt,
  before skip=10pt, after skip=10pt,
  #1
}

% ── Editorial left-rule — the PDF twin of .di-editorial ───────────────────
\newtcolorbox{editorial}[1][]{
  enhanced, breakable,
  colback=SurfacePage, colframe=EditorialRule,
  boxrule=0pt, leftrule=2pt, arc=0pt,
  left=12pt, right=0pt, top=4pt, bottom=4pt,
  before skip=10pt, after skip=10pt,
  #1
}

% ── Code blocks — match the HTML `pre` rule (8pt padding, hairline, r=2pt) ─
\ifdefined\Shaded
  \renewenvironment{Shaded}{%
    \begin{tcolorbox}[
      enhanced, breakable,
      colback=SurfaceCode, colframe=BorderHairline,
      boxrule=0.4pt, arc=2pt,
      left=8pt, right=8pt, top=8pt, bottom=8pt,
      before skip=8pt, after skip=8pt
    ]%
  }{%
    \end{tcolorbox}%
  }
\fi

% ── Block quotes — match the HTML blockquote rule ─────────────────────────
\renewenvironment{quote}{%
  \begin{tcolorbox}[
    enhanced, breakable,
    colback=SurfaceSubtle, colframe=InteractiveAccent,
    boxrule=0pt, leftrule=3pt, arc=2pt,
    left=9pt, right=9pt, top=6pt, bottom=6pt,
    before skip=8pt, after skip=8pt
  ]\itshape\color{ContentSecondary}%
}{%
  \end{tcolorbox}%
}

% ── Captions, lists, tables ───────────────────────────────────────────────
\captionsetup{font=small, labelfont={bf,color=ContentPrimary},
              textfont={color=ContentMuted},
              justification=justified, skip=6pt}

\setlist[itemize]{leftmargin=*, itemsep=0.3em, parsep=0.15em, topsep=0.4em}
\setlist[enumerate]{leftmargin=*, itemsep=0.3em, parsep=0.15em, topsep=0.4em}

\renewcommand{\arraystretch}{1.25}

% ── Paragraph shape ───────────────────────────────────────────────────────
\setlength{\parskip}{0.5\baselineskip}
\setlength{\parindent}{0pt}
\setlength{\emergencystretch}{3em}

% Quarto's callouts are plain tcolorbox environments. By default tcolorbox
% typesets its contents in a parbox, which zeroes \parskip — so paragraphs
% inside a callout run together while the surrounding document separates them.
% parbox=false restores normal paragraph handling inside every box.
\tcbset{parbox=false}

% ── Stat tiles — the PDF twin of .di-stats / .di-stat ─────────────────────
% HTML lays these out as a grid; print gets a rule-separated row of figures.
\newtcolorbox{statrow}[1][]{
  enhanced, breakable,
  colback=SurfacePage, colframe=SurfacePage,
  boxrule=0pt, arc=0pt,
  left=0pt, right=0pt, top=4pt, bottom=4pt,
  before skip=10pt, after skip=10pt,
  #1
}
\newcommand{\statvalue}[1]{{\HeadingFont\fontsize{20pt}{22pt}\selectfont\color{ContentPrimary}#1}}
\newcommand{\statlabel}[1]{{\footnotesize\color{ContentMuted}#1}}
